\section{Discussion}
\label{sec:discussion}

The yield of \SI{71.4}{\percent} (Equation~\ref{eq:yield}) is typical for this
experiment. Losses occurred mainly during recrystallisation, where some product
remained dissolved in the mother liquor. The narrow melting range and the single
spot on the TLC plate (Figure~\ref{fig:tlc}) indicate a pure product, which the
negative ferric chloride test confirms.

The crude product contained some unreacted salicylic acid, as shown by its
higher mass in Table~\ref{tab:yield} and a faint violet colour with
\ce{FeCl3}. Recrystallisation removed this impurity. Acetylation is effectively
irreversible under these conditions because the acetate leaving group,
\ce{CH3COO-}, is protonated in the acidic medium~\cite{clayden2012}.

\subsection{Conclusion}
Aspirin was prepared in good yield and high purity. The reaction mechanism and
the role of the acid catalyst were discussed in Section~\ref{sec:intro}.
